\scanpage{118}
\mspage{33}
% archive sheet number: 117

% L01
\[
H^*(U(m),\Lambda)\unsure{{-}|\longrightarrow}H^*(U(m'),\Lambda),\quad\text{pour }m\mid m'
\]

% L02-04
est justiciable de 4.6 (où on \deleted{\unsure{prendra}}\inserted{\unsure{remplacera|remplace}} $n$ par $m'/m$), donc est nul \inserted{en degrés $\ne0$}, si $nm\mid m'$, \inserted{sous réserve du th. de pureté \unsure{coh.|ch.}} \deleted{D'ailleurs,} Cela prouve que le deuxième \unsure{mbr} de $(*)$ est nul en dim $>0$, et en dim. 0 il est évidemment isomorphe à $\Lambda$. Donc on trouve

% L05-07
\uline{Corollaire} 4.9. Dans les notations de 4.8, et supposant que $X$ et les $X(m)$ satisfont au th. de pureté coh., \deleted{\unsure{on dit}} on a

% L08-09
\[
\left\{\begin{aligned}
H^i(U(\infty),\Lambda)&=0\quad\text{si }\deleted{\unsure{\Lambda}}\ i\ne0\\
H^0(U(\infty),\Lambda)&\simeq\Lambda
\end{aligned}\right.
\tag*{(4.9.1)}
\]

% L10-14
\deleted{\uline{Remarques} 4.10.\unsure{a)} Si on suppose seulement que les $X(m)$ satisfont au th. de pureté \inserted{coh.} en dim $\leq N$ (condition vide si $N=1$ !), alors on trouve que (4.9.1) est encore vrai si $i\leq N$.}

% L15-17
\deleted{b) On peut exprimer encore de façon imagée le résultat de 4.9 en disant que \inserted{\unsure{à $p$-torsion près}} \unsure{cohomologiquement} $U$ se comporte comme un}

% L18-19
\uline{Corollaire} 4.10. \inserted{Sous les conditions de 4.9,} Soit $\Gamma$ un groupe profini, considérons un homomorphisme de groupes profinis

% L20
\[
\unsure{\varphi|\phi}:G\deleted{(\unsure{\infty})}\longrightarrow\Gamma,
\tag*{(4.10.1)}
\]

% L21
et soit $H$ le noyau de cet homomorphisme.
